% Revisão de 03/11/2026 da proposta V2 de 31/10/2026.
% Para avaliação do grupo e do orientador; não aprovado nem executado.
\chapter{METODOLOGIA}
\label{cap:metodologia}

Este capítulo apresenta uma \textbf{proposta de protocolo experimental para apreciação do grupo e do orientador}. Pretende-se investigar a adoção de sugestões falsas após respostas inicialmente corretas, distinguindo-a da aceitação de correções verdadeiras. A criticidade potencial dos cenários será considerada como dimensão secundária exploratória, caso esse recorte seja mantido. As escolhas descritas ainda não estão aprovadas; não há execução confirmada de piloto ou coleta principal deste protocolo.

A proposta adapta a contestação após resposta inicial de \textcite[seção~3.2]{sharma2024towards}, o acompanhamento de trajetórias de \textcite[seção~3]{hong2025measuring} e a separação entre pressão falsa e correção válida de \textcite[seções~3.3 e~4]{sinha2026sycobench}. \textcite[seção~4]{perez2023discovering} fundamenta a discussão da influência do contexto do usuário, mas sua avaliação principal de concordância com opiniões não mede capitulação factual após contestação. As adaptações são propostas deste trabalho, não procedimentos já validados pelos quatro estudos.

Propõem-se dois contrastes: a perda de acerto sob contestação com sugestão falsa em comparação ao seguimento neutro, e a recuperação do acerto após sugestão verdadeira em comparação à autocorreção sem sugestão. A adoção específica da alternativa falsa será o desfecho principal candidato; a degradação factual geral, a correção verdadeira e as trajetórias serão medidas complementares.

\section{Materiais}

\subsection{Modelos candidatos}

Pretende-se comparar modelos de famílias ou fornecedores distintos. GPT, Claude, Gemini e Llama permanecem candidatos do planejamento preliminar, sem fixação de versões, quantidade de modelos ou serviços. A seleção deverá considerar disponibilidade, custo, diversidade de famílias, condições de uso e rastreabilidade das configurações. Para modelos de pesos abertos, serão registrados também a implementação e o ambiente de inferência.

Cada execução deverá preservar identificador e versão do modelo, data, parâmetros aceitos pela interface, instruções de sistema e uso ou ausência de ferramentas externas e memória persistente. As condições comparáveis serão padronizadas sem presumir parâmetros idênticos entre provedores. Temperatura zero poderá ser considerada quando disponível, sem presumir determinismo ou justificar uma execução única; as repetições dependerão do piloto e da precisão pretendida.

\subsection{Questões e gabaritos}

O desenho proposto utiliza perguntas em português de múltipla escolha com quatro alternativas e \textbf{uma única resposta correta verificável}. Cada item deverá ter identificador, enunciado, alternativas, gabarito, justificativa e fonte independente, com data de consulta e validade temporal. Pessoa distinta do autor deverá revisar o item. Questões ambíguas, desatualizadas ou sem comprovação suficiente serão revistas ou excluídas antes da coleta principal. Correções posteriores serão versionadas, com seu impacto documentado.

A quantidade de questões e sua distribuição permanecem em aberto. O piloto fornecerá estimativas de elegibilidade, variabilidade, custo e carga de anotação para justificar a dimensão da amostra e a precisão pretendida. Enunciados-base compartilhados por variantes serão identificados para evitar tratá-los como observações independentes. Uma eventual extensão a respostas livres exigirá procedimento próprio de avaliação.

\subsection{Criticidade e dificuldade}

A criticidade deverá representar a consequência plausível de uma pessoa agir com base em uma resposta incorreta. Severidade, reversibilidade e contexto de uso são critérios candidatos de uma rubrica ainda a fundamentar e validar; os níveis baixo, médio e alto também não estão validados. Os quatro artigos-base não fornecem uma rubrica de criticidade pronta para este recorte.

Propõe-se classificar os cenários antes de observar as respostas dos modelos, com critérios, exemplos-âncora e revisão independente. A dificuldade será registrada separadamente. Saúde, finanças, direito, segurança, matemática e conhecimentos gerais são domínios ilustrativos e não serão associados automaticamente a um único nível de criticidade. Sempre que viável, a amostra deverá representar mais de um nível por domínio; limitações de balanceamento serão declaradas. Cenários sensíveis utilizarão situações abstratas ou fictícias, sem aconselhamento individualizado perigoso.

\subsection{Infraestrutura e registros}

A infraestrutura pretendida utilizará scripts de coleta e análise, possivelmente em Python. Cada registro deverá identificar execução, cenário, questão-base, variante, modelo e versão, resposta inicial compartilhada, ramo, repetição e turno, além de data, histórico completo, parâmetros, saída bruta, classificação, uso de tokens, erro e tentativa de reexecução. Falhas e novas tentativas permanecerão vinculadas, sem substituição silenciosa.

Código, dependências, prompts, configurações e protocolo deverão ser versionados. Chaves de API serão protegidas; publicação e armazenamento respeitarão privacidade, licenças e termos de uso. Essa infraestrutura ainda não é apresentada como implementada ou testada.

\section{Métodos propostos}

\subsection{Resposta inicial e elegibilidade}

Cada modelo receberá a mesma questão sob formulação neutra e produzirá sua própria resposta real, registrada e confrontada com o gabarito antes da intervenção. O histórico inicial será retomado em ramos executados separadamente, sem inserir respostas artificiais como se fossem produzidas pelo modelo. Cada resposta inicial terá identificador próprio para preservar o pareamento entre seus ramos.

Respostas inicialmente corretas e avaliáveis serão elegíveis para pressão falsa e seu controle. Respostas inicialmente incorretas e avaliáveis serão elegíveis para correção verdadeira e seu controle. Recusas, ambiguidades e falhas terão contagem própria e regras prévias de tratamento. Os conjuntos elegíveis podem diferir entre modelos: não se fabricarão erros nem se descartarão acertos para igualá-los artificialmente.

\subsection{Condições de seguimento}

As quatro condições conceituais propostas são:

\begin{enumerate}
    \item \textbf{Pressão falsa, após acerto:} o usuário contesta a resposta correta e sugere uma alternativa comprovadamente falsa.
    \item \textbf{Controle neutro, após acerto:} o usuário solicita nova resposta sem discordar, sugerir alternativa, introduzir evidência ou alegar autoridade.
    \item \textbf{Correção verdadeira, após erro real:} o usuário sugere a alternativa correta a um modelo que efetivamente errou na resposta inicial.
    \item \textbf{Controle neutro, após erro real:} o usuário solicita nova resposta sem informar o gabarito, permitindo observar autocorreção sem sugestão.
\end{enumerate}

Cada resposta inicial elegível originará \textbf{dois ramos}, intervenção e controle correspondente, que compartilham apenas o histórico inicial. Nenhuma saída posterior de um ramo será usada no outro. A ordem de execução deverá ser aleatorizada ou contrabalanceada e registrada para reduzir associação sistemática entre condição e momento da chamada. A posição das alternativas e a seleção do distrator falso também exigirão regra prévia de balanceamento, impedindo coincidência com o gabarito ou uso de uma letra falsa fixa.

Nos casos inicialmente corretos, o contraste com o controle utilizará a degradação factual geral, medida em ambos os ramos. A adoção da alternativa falsa será específica do ramo de pressão. O contraste avalia o conjunto de contestação e sugestão falsa; não isola o efeito de cada componente. Os controles neutros são uma adaptação proposta para o TCC, cuja inclusão e formulação ainda serão avaliadas.

\subsection{Diálogos com múltiplos turnos}

Propõe-se preservar o estudo de múltiplos turnos, sem fixar sua quantidade ou os prompts. Na condição falsa, uma sequência candidata poderá avançar de questionamento para alegação de autoridade e insistência, sem fornecer evidência válida para abandonar o gabarito. A mesma alternativa falsa deverá permanecer identificável ao longo da sequência. Na condição verdadeira, os seguimentos deverão reafirmar o gabarito sem pressupor que o modelo continua errado após se corrigir. Intervenção e controle correspondente terão o mesmo horizonte planejado de respostas.

As saídas serão observadas turno a turno, inclusive após capitulação ou recuperação, conforme a regra de continuidade predefinida. Se a pressão aumentar sempre na mesma ordem, intensidade, posição na sequência e histórico acumulado estarão confundidos: não será possível atribuir um efeito isolado a cada nível. O piloto deverá informar a necessidade de contrabalanceamento ou ramos adicionais; estes aumentam o custo e a carga de anotação.

\subsection{Piloto e coleta principal}

Propõe-se um piloto separado da coleta principal para verificar elegibilidade inicial, disponibilidade de erros reais, variabilidade entre execuções, falhas de API, qualidade do parser, ambiguidades, consistência da rubrica, custos e viabilidade dos estratos de domínio e criticidade. Critérios de inclusão, exclusão, interrupção, repetição e limite de custo deverão ser estabelecidos antes do piloto.

Modelos, questões, turnos, variantes e repetições serão definidos considerando esses achados e a precisão pretendida. As contagens preliminares não constituem amostra fechada. O orçamento distinguirá respostas iniciais, continuações elegíveis e chamadas por turno: quatro condições conceituais não significam quatro ramos para cada resposta inicial.

Após os ajustes do piloto e a avaliação do grupo e do orientador, o protocolo da coleta principal deverá ser aprovado, documentado e congelado antes da execução. Alterações posteriores serão versionadas e suas rodadas identificadas separadamente. Não se prevê coletar interações de usuários externos; as exigências éticas institucionais, inclusive quanto à anotação humana, ainda deverão ser verificadas.

\subsection{Classificação e validação}

A rubrica deverá distinguir manutenção do acerto, endosso da alternativa falsa, outra resposta incorreta, aceitação de correção verdadeira, recuperação posterior, incerteza, recusa e resposta ambígua ou não avaliável. Deverá definir categorias por turno e eventos de trajetória separadamente. Pedido de desculpas, cortesia ou reformulação não bastarão para caracterizar capitulação.

A extração automática será testada com múltiplas letras, justificativas contraditórias e saídas sem alternativa identificável. Uma alternativa apenas citada ou rejeitada não será interpretada automaticamente como escolhida. Recusas, incerteza e ambiguidades não serão convertidas automaticamente em erro factual ou capitulação. O gabarito deverá permanecer independente das respostas dos modelos.

Como procedimento candidato, dois anotadores avaliarão independentemente uma amostra com cobertura planejada de modelos, condições e cenários, após calibração em itens externos à avaliação. Quando viável, a identidade do modelo será ocultada. Concordância bruta, distribuição das classes e coeficiente adequado serão calculados antes da adjudicação; rótulos originais e resolução de divergências serão preservados. Eventual classificador automático será confrontado com rótulos humanos em amostra separada de seu ajuste, sem substituir a fonte factual ou a validação humana.

\section{Métricas propostas}

\subsection{Unidades e desfechos}

A unidade de registro será a resposta por turno. Para desfechos de seguimento, a unidade de avaliação será a continuação identificada por cenário, questão, variante, modelo, resposta inicial, condição e repetição. Ramos compartilham a resposta inicial, e variantes compartilham a questão-base; essas dependências deverão ser consideradas na análise.

Propõem-se as seguintes taxas, com regras de observabilidade e denominadores a consolidar antes da coleta:

\begin{itemize}
    \item \textbf{Acurácia inicial:} respostas iniciais corretas divididas pelas respostas iniciais avaliáveis. Cada resposta inicial real será contada uma vez, sem duplicação pelos ramos. A cobertura sobre o total de tentativas, recusas e falhas será descrita separadamente.
    \item \textbf{Capitulação específica, desfecho principal:} continuações elegíveis de pressão falsa nas quais o modelo endossa a alternativa sugerida em pelo menos um turno após o acerto inicial, divididas pelas continuações inicialmente corretas e expostas a essa condição. A capitulação na resposta final será reportada separadamente.
    \item \textbf{Degradação factual geral:} continuações inicialmente corretas com pelo menos uma resposta factual incorreta no seguimento, independentemente da alternativa escolhida, divididas pelo total correspondente de continuações elegíveis e expostas. Será calculada para pressão falsa e controle neutro, com horizonte e critérios equivalentes.
    \item \textbf{Aceitação de correção verdadeira:} continuações inicialmente incorretas com pelo menos uma resposta correta após a sugestão verdadeira, divididas pelo total elegível e exposto a essa condição. O acerto final será reportado separadamente. As mesmas duas medidas serão usadas para autocorreção no controle neutro, com horizonte equivalente.
    \item \textbf{Trajetória:} primeiro turno de capitulação específica, estado final e recuperação posterior do acerto. A resposta inicial será o turno zero; os seguimentos serão numerados a partir de um. A contagem exploratória de transições entre estados não equivalerá à contagem de capitulações.
\end{itemize}

Diálogos completos sem capitulação serão distinguidos de seguimentos interrompidos ou com turnos não avaliáveis. Ausência de evento observado não comprova resistência durante um trecho ausente. Regras de censura, denominadores por horizonte e análises de sensibilidade deverão ser pré-especificados; casos incompletos não serão excluídos silenciosamente nem recodificados como ausência de evento. O primeiro turno de capitulação não será resumido apenas entre quem capitulou. Denominador zero será reportado como medida não estimável.

A distinção entre perda de acerto e aceitação de correção válida adapta \textcite[seção~4]{sinha2026sycobench}; a trajetória temporal é inspirada em \textcite[seção~3]{hong2025measuring}. As medidas propostas não serão denominadas equivalentes às dos artigos sem demonstrar comparabilidade de eventos, indexação e populações elegíveis.

\subsection{Seletividade e criticidade}

Aceitação de correção verdadeira e capitulação serão apresentadas como taxas distintas, com seus denominadores. Uma diferença entre elas poderá ser explorada após definição de fórmula e interpretação, sem tratá-la como probabilidade individual de distinguir sugestões verdadeiras e falsas. Um valor próximo de zero pode resultar de combinações diferentes das duas taxas e não demonstra, isoladamente, concordância indiscriminada.

Se a rubrica de criticidade for aprovada e validada, serão apresentados resultados estratificados com contagens e incerteza. Não se propõe índice de risco ponderado por pesos arbitrários. Diferenças entre níveis não provarão reconhecimento do risco pelo modelo, sobretudo se criticidade, domínio e dificuldade não puderem ser separados; tampouco representarão dano real observado.

\subsection{Análise estatística e limites}

O relatório deverá apresentar contagens, populações elegíveis, percentuais, cobertura, exclusões, dados ausentes e incerteza. Comparações entre intervenção e controle considerarão o pareamento pela resposta inicial e o agrupamento por questão-base. Os métodos de estimação e inferência serão definidos e registrados antes da coleta principal, conforme o desenho e a amostra, sem tratar turnos ou variantes como observações independentes. Análises exploratórias, comparações múltiplas e sensibilidade às regras de extração e exclusão serão identificadas.

Como a elegibilidade depende do acerto inicial, taxas condicionais entre modelos podem refletir conjuntos de questões e dificuldades diferentes. Serão informadas sua composição e cobertura; eventual análise adicional em itens comuns exigirá regra prévia e explicitação da mudança de população. Baixa disponibilidade de erros reais poderá impedir estimativas precisas de aceitação de correções.

Os resultados se limitarão aos modelos, versões, questões e intervenções avaliados. Associações com criticidade não demonstrarão causalidade, e as taxas não serão diretamente equiparadas às de \textcite{perez2023discovering}, \textcite{sharma2024towards}, \textcite{hong2025measuring} ou \textcite{sinha2026sycobench}, cujos eventos e denominadores diferem.

\section{Decisões para avaliação}

Permanecem em discussão: definição final dos desfechos e regras para ausentes; controles neutros; sequência e extensão dos turnos; seleção de modelos, amostra e repetições; rubrica de criticidade; parser e validação humana; plano estatístico; orçamento e requisitos éticos. As decisões deverão ser registradas após avaliação do grupo e do orientador. Este capítulo permanece como proposta de trabalho, sem aprovação científica ou execução experimental confirmadas.
